% 한글 독립 문서. 권장 컴파일러: XeLaTeX.
\documentclass[10pt,a4paper]{article}
\usepackage[margin=20mm]{geometry}
\usepackage{kotex}
\usepackage{amsmath,amssymb}
\usepackage{array,longtable,booktabs,calc}
\usepackage{xurl,xcolor}
\usepackage[unicode,hidelinks]{hyperref}
\setlength{\parindent}{0pt}
\setlength{\parskip}{5pt}
\setlength{\emergencystretch}{3em}
\setlength{\tabcolsep}{4pt}
\renewcommand{\arraystretch}{1.15}
\newcolumntype{L}[1]{>{\raggedright\arraybackslash}p{#1}}
\newcolumntype{C}[1]{>{\centering\arraybackslash}p{#1}}
\newcommand{\yes}{\ensuremath{\circ}}
\newcommand{\maybe}{\ensuremath{\triangle}}
\newcommand{\no}{\ensuremath{\times}}
\newcommand{\yin}{\href{https://doi.org/10.1007/s11116-022-10319-3}{Yin and Cherchi (2024)}}
\newcommand{\loa}{\href{https://doi.org/10.1007/s11116-026-10751-9}{Loa, Lee, and Circella (2026)}}
\newcommand{\elk}{\href{https://arxiv.org/abs/2406.11590}{Elkhouly and Alhadidi}}
\newcommand{\zhang}{\href{https://arxiv.org/abs/2209.07980}{Zhang et al.}}
\newcommand{\wei}{\href{https://www.nature.com/articles/s41599-024-03271-3}{Wei et al. (2024)}}
\newcommand{\experience}{\href{https://doi.org/10.1177/03611981231155181}{Yin, Cherchi, and Nettle (2023)}}
\newcommand{\cats}{\href{https://doi.org/10.1371/journal.pone.0262496}{Cats et al.}}
\newcommand{\rapid}{\href{https://doi.org/10.1177/03611981221093632}{Etminani-Ghasrodashti et al. (2022)}}
\newlength{\fourwidth}
\setlength{\fourwidth}{\textwidth-6\tabcolsep}
\newlength{\applywidth}
\setlength{\applywidth}{\textwidth-14\tabcolsep}
\begin{document}
\begin{center}
{\Large\bfseries $p$의 driving factor 선택과정 정리}
\end{center}

%----------------------------------
\section{요약}
이 문서는 \textbf{ER-DRO 모형에서 불확실한 로봇택시 선택비율 $p$를 추정할 변수}를 고른다. 순서는 \textbf{후보군 $\rightarrow$ 기준 $\rightarrow$ 기준 적용 $\rightarrow$ 선택} 결과이다. 
\[
 p_c\text{의 관측값} =y_c=\frac{\text{도시 }c\text{의 로봇택시 통행}}{M_c}.
\]
분모 $M_c$와 분자는 같은 지리범위ㆍ기간ㆍ통행 단위로 맞춘다. 이용자료가 완료 통행만 포함하면 $y_c$는 공급조건이 반영된 실현 이용비율이다.

변수 선택은 마지막 선택표에서 외생과 내생(운영 결정에 의존)을 구분한다. 결정에 의존하는 효과를 최적화에 반영할 때는 결정의존 모형으로 연결해야 한다.

향후 $p=f(x)+Q(x)\epsilon$ 형태와 결정의존 확장도 고려한다. 이 문서에서는 확장에 필요한 조건만 선택기준에 반영한다.

\textbf{결론.} \textcolor{blue}{\textbf{1차 선택 결과: 외생변수로는 ''무차량 가구 비율'', ''연령 구조'', 내생변수로는 ''상대 요금'', ''상대 대기시간'', ''차내 지원기능'' 선택한다.}}

%------------------------------------
\section{이미 정해놓은 후보군}
\subsection{Driving factors of $p$}
아래 표는 driving factor 후보 23개이며, \texttt{A.Background\_Research}에서 정리한 driving factor들에서 가져왔다. 도시 수준 $p$의 효과가 모두 입증됐다는 뜻은 아니다.

%-------A. 기존 $p$ 후보 8개-----------
{\footnotesize
\begin{longtable}{@{}L{\fourwidth*\real{0.16}}L{\fourwidth*\real{0.30}}L{\fourwidth*\real{0.34}}L{\fourwidth*\real{0.20}}@{}}
\toprule 변수 & 근거 논문 & 논문에서 확인된 범위ㆍ결과 & 측정자료\\
\midrule\endfirsthead
\toprule 변수 & 근거 논문 & 논문에서 확인된 범위ㆍ결과 & 측정자료\\
\midrule\endhead
\bottomrule\endlastfoot
\textbf{상대 요금} & \yin, \emph{Willingness to pay for automated taxis} & 일반택시와 자동택시의 선택실험. 요금에 따른 선택 차이를 분석. 자사 요금은 운영자 결정. & 양 서비스의 제시요금 + 선택결과\\
\textbf{상대 대기시간} & \yin & 대기시간에 따른 선택 차이를 분석. 자사 대기시간은 차량 수ㆍ배치에 의존. & 양 서비스의 대기시간 + 선택결과\\
\textbf{차내 통행시간} & \yin & 요금ㆍ통행시간의 선택상충 관계 분석. & 양 서비스의 제시 통행시간\\
\textbf{신뢰ㆍ가격가치ㆍ이용습관} & \wei, \emph{Using the Extended UTAUT}; \experience & 수용의향ㆍ사용행동의 설문 근거. Wei에는 신뢰 변수가 없음. 태도 점수 자체는 선택확률이 아님. & 개인 태도ㆍ습관 설문\\
\textbf{연령 및 로봇택시 관련 지식} & \yin; \wei & 단순한 연령 효과는 확정하기 어려움. 일반 AV 지식과 실제 서비스 운영 인지를 구분. & 개인 특성과 선택자료의 연결\\
\textbf{사용자 리뷰 및 차내 지원기능} & \yin & 좋은 리뷰, 목적지 변경, 운영자 소통 등 서비스 속성의 선택효과 분석. & 선택실험의 서비스 속성\\
\textbf{로봇택시 이용 경험} & \experience, \emph{Living in a City Where Automated Taxis are Operating} & 운영도시 거주ㆍ개인 이용 경험과 선호의 관계. 경험에 따른 단순한 증가 효과는 확정되지 않음. & 개인 이용 경험; 도시 운영ㆍ개시 기록\\
\textbf{대중교통 등 대안의 경쟁력} & \cats, \emph{Beyond the dichotomy} & 일반 ride-hailing 근거. 통행시간 경쟁력에 따라 대체ㆍ보완이 공존. & 대안 가용성ㆍ요금ㆍ시간; GTFS, NTD\\

%--------B. $M$ 표에서 이동한 후보 7개------

\textbf{무차량 가구 비율} & \elk; \loa & Chicago 구역의 일반 호출 이용과 양(+)의 연관. 첨부 문서는 SF 로봇택시 통행 발생에서도 양(+)의 연관을 보고. & ACS B08201\\
\textbf{소득} & \elk; \loa; \yin; \wei & ChicagoㆍSF 집계자료에서는 양(+). 선택실험ㆍ설문에서는 같은 효과가 일관되게 확인되지 않음. & ACS B19013\\
\textbf{고용 및 경제활동} & \elk; \loa & 경제활동 비율은 Chicago 일반 호출 이용, 고용밀도는 SF 로봇택시 통행 도착과 양(+)의 연관. & LEHD/LODES; ACS B23025\\
\textbf{인구밀도} & \elk; \zhang; \loa & 일반 호출에서는 공간적 차이, SF 로봇택시 이용량에서는 음(-)의 연관. & ACS 인구 + 면적\\
\textbf{토지이용 혼합ㆍ접근성ㆍ상업활동} & \zhang; \loa & Chicago에서는 비선형ㆍ공간 이질적 관계. SF에서는 토지이용 혼합과 양(+)의 연관. & EPA SLD\\
\textbf{대중교통ㆍ보행환경} & \zhang; \cats; \loa & 일반 호출에서 대체ㆍ보완이 공존. 첨부 문서는 SF에서 빈도ㆍ보행성과 양(+)의 연관을 보고. & EPA SLD; GTFS; NTD\\
\textbf{식당ㆍ여가시설 밀도} & \loa & 첨부 문서는 SF 로봇택시 통행과 양(+)의 연관을 보고. $M$과 $p$의 효과를 분리한 결과는 아님. & Census CBP; OpenStreetMap\\


%-----C. 이전 대화의 추가 후보 8개--------

\textbf{연령 구조} & \loa; \rapid & 첨부 문서ㆍ대화의 젊은 층 이용 근거. 개인 선택실험에서는 연령의 단순 부호를 확정하기 어려움. & ACS B01001\\
\textbf{학생 비율} & \rapid & Arlington 캠퍼스 파일럿 이용 근거. 학생의 무료 이용 지원 조건을 고려해야 함. & ACS 재학자료; 지역 학생인구\\
\textbf{교육ㆍ기술 친화도} & 이전 대화의 SAV 선택실험ㆍ리뷰 문헌 & 추가 후보. 학력과 기술 친화도는 다른 개념이며, 세부 효과는 확인 필요. & ACS 학력; 기술 친화도 설문\\
\textbf{기존 ride-hailing 습관} & \wei; 기존 일반 호출 연구 & 습관ㆍ과거 사용이 채택과 연결될 가능성. 과거 이용량은 습관의 대리자료. & 공통 설문; 과거 일반 호출 이용자료\\
\textbf{안전 인식ㆍ노출} & 이전 대화의 SAV 연구; Khan(2022)은 서지ㆍ결과 재확인 필요 & 안전 인식 후보. 인지된 안전과 사고ㆍ범죄 노출은 구분해야 함. & 안전 설문; 사고ㆍ노출자료\\
\textbf{규제 허용 여부} & 수요효과 논문이 아니라 실행가능성 자료 & 허용 조건을 충족하지 않으면 운행 불가. 선호효과와 별도로 처리. & 관할 기관의 허가ㆍ운행 조건; NCSL\\
\textbf{기후ㆍODD} & 이전 대화의 운행 가능 조건 & 기상과 운행 가능 범위의 적합성. 선호와 운행 가용성을 구분. & NOAA 기상; 운영자 ODD\\
\textbf{서비스 개시 후 경과기간} & \experience; 개시ㆍ운행 기록 & 인지ㆍ경험의 누적을 나타내는 후보. 기간 자체의 효과는 별도 검증 필요. & 회사 발표; 허가ㆍ개시 기록\\
\end{longtable}}

%-----------------------------------
\section{$p$의 driving factor 선택기준}
{\small
\begin{longtable}{@{}L{0.07\textwidth}L{0.72\textwidth}L{0.15\textwidth}@{}}
\toprule 번호 & 기준 & 성격\\
\midrule\endfirsthead
\toprule 번호 & 기준 & 성격\\
\midrule\endhead
\bottomrule\endlastfoot
0 & 종속변수의 분모ㆍ지리범위ㆍ기간ㆍ통행 단위가 일치 & 사전\\
1 & $p$와 관련된 문헌 근거와 해석 가능한 이유가 있음 & 사전+사후\\
2 & 모든 후보 도시에서 같은 정의ㆍ출처로 결정시점에 측정 가능 & 사전\\
3 & 해석 가능; 중복 변수 최소화; 표본 수에 맞춰 변수 수 제한 & 사전\\
4 & 도시 규모보다 비율ㆍ밀도ㆍ상대 서비스 조건을 측정 & 사전\\
5 & 운영 결정에 의존하는지 구분하고 모형에서의 역할을 정함 & 사전\\
6 & 후보 도시 값이 학습자료의 범위와 서비스 조건에 부합 & 사전+사후\\
7 & 적합 후 잔차에 공변량에 따른 설명되지 않은 패턴이 없음 & 사후\\
8 & 결정의존 확장 시 결정 변화에 따른 $p$의 반응을 추정할 자료가 있음 & 확장 시\\
\end{longtable}}

\subsection{각 기준에 대한 검토}
\begin{enumerate}
\setcounter{enumi}{-1}
\item \textbf{종속변수 일치.} 로봇택시 통행/$M$의 분자ㆍ분모를 맞춘다. 택시 두 종류 사이의 선택확률은 전체 이동기회에서의 $p$와 다르다.
\item \textbf{문헌 근거.} 실제 로봇택시, SAV 파일럿, 일반 호출의 근거를 구분한다. 부호가 다르면 서비스ㆍ지역 차이를 확인하고 보류한다.
\item \textbf{공통 측정.} 같은 도시 경계ㆍ연도ㆍ변수 정의를 사용한다. 동일 설문이 없는 신뢰ㆍ안전 인식 등은 자료 확보 전 보류한다.
\item \textbf{간결성.} 비슷한 개념은 대표변수 하나부터 선택한다. 독립 도시 수를 고려하며, $n/10$은 보수적인 참고 경험칙으로만 사용한다.
\item \textbf{변수 형태.} 인구는 $M$의 규모인자로 둔다. $p$에는 무차량 가구 비율ㆍ시설 밀도ㆍ상대 요금ㆍ상대 시간 등을 사용한다.
\item \textbf{외생ㆍ내생 구분.} 도시 특성은 외생 후보, 운영자가 정하거나 결정에 따라 달라지는 서비스 조건은 내생 후보로 표시한다. 내생이라는 이유만으로 후보에서 제외하지 않는다.
\item \textbf{도시 전이.} 구역 수준 또는 다른 국가의 결과가 후보 도시에 적용되는지 확인한다. 학습범위 밖의 값은 외삽으로 표시한다.
\item \textbf{잔차 확인.} $f(x)+\epsilon$이면 잔차를, $f(x)+Q(x)\epsilon$이면 $Q(x)$로 나눈 표준화 잔차를 확인한다. 평균ㆍ퍼짐 외에 분포 패턴도 점검한다.
\item \textbf{결정의존 확장.} 요금ㆍ차량 수 등의 변화와 그때의 선택ㆍ서비스 결과가 함께 필요하다. 대기시간이 수요에도 반응하면 그 관계를 함께 고려한다.
\end{enumerate}

\section{후보군에 기준 적용: $p$}
첨부 문서의 $M$ 적용표와 같은 \textbf{변수 / 1 / 2 / 3 / 4 / 5 / 판정 / 이유ㆍ쓰는 형태} 형식이다. 기준 0은 공통 전제이며, 6ㆍ7ㆍ8은 데이터ㆍ적합 결과가 있어야 판정하므로 추후 확인한다.

\yes: 설계상 적용 가능, \maybe: 추가 확인, --: 해당 없음. 기준 1의 +ㆍ-는 문헌의 연관 방향이다. 기준 5의 외ㆍ내ㆍ혼ㆍ별은 각각 외생ㆍ내생ㆍ구성요소에 따라 다름ㆍ별도 조건을 뜻한다. \textbf{선택은 1차 선택이며 자료 검증 완료를 뜻하지 않는다.}


{\footnotesize
\begin{longtable}{@{}L{\applywidth*\real{0.22}}*{5}{C{\applywidth*\real{0.04}}}C{\applywidth*\real{0.08}}L{\applywidth*\real{0.50}}@{}}
\toprule 변수 & 1 & 2 & 3 & 4 & 5 & 판정 & 이유 / 쓰는 형태\\
\midrule\endfirsthead
\toprule 변수 & 1 & 2 & 3 & 4 & 5 & 판정 & 이유 / 쓰는 형태\\
\midrule\endhead
\bottomrule\endlastfoot
상대 요금 & - & \maybe & \yes & \yes & 혼 & 선택 & 자사 요금은 내생, 경쟁요금은 외생 후보. 동일 통행의 요금 차이로 사용하며 결정의존 확장에 반영.\\
상대 대기시간 & - & \maybe & \yes & \yes & 혼 & 선택 & 자사 대기는 차량 수ㆍ배치에 의존. 경쟁서비스 대기와 비교하며 결정의존 확장에 반영.\\
차내 통행시간 & - & \maybe & \yes & \yes & 혼 & 보류 & 교통조건은 외생, 경로ㆍ운영정책이 바꾸는 부분은 내생. 상대 통행시간 자료 확인 필요.\\
신뢰ㆍ가격가치ㆍ이용습관 & \maybe & \maybe & \maybe & \yes & 혼 & 보류 & 공통 설문 필요. 서로 다른 개념을 한 변수로 합치지 않음.\\
연령ㆍ관련 지식 & \maybe & \maybe & \maybe & \yes & 외 & 일부 & 연령은 C의 연령 구조로 통합. 운영 인지ㆍ지식은 설문자료 확인 전 보류.\\
리뷰ㆍ차내 지원기능 & + & \maybe & \maybe & \yes & 혼 & 일부 & 지원기능은 내생 후보로 선택. 리뷰는 과거 사용의 결과이므로 별도 보류.\\
로봇택시 이용 경험 & \maybe & \maybe & \yes & \yes & 혼 & 보류 & 현재 결정 이전의 경험인지 확인. 개인 경험을 도시 운행기간으로 바로 대체하지 않음.\\
대안의 경쟁력 & \maybe & \maybe & \maybe & \yes & 외 & 보류 & 대안 요금ㆍ시간ㆍ가용성을 비교. 대중교통ㆍ보행환경과 중복 확인.\\

%-----M에서 이동한 후보----------------
무차량 가구 비율 & + & \yes & \yes & \yes & 외 & 선택 & 차량을 보유하지 않은 가구 비율. 자가용 대안 제약을 직접 설명하며 ACS로 공통 측정 가능.\\
소득 & \maybe & \yes & \maybe & \yes & 외 & 보류 & 집계 이용과 선택실험의 효과가 일관되지 않음. 소득 또는 경제활동 비율을 하나씩 추가해 비교.\\
고용ㆍ경제활동 & + & \yes & \maybe & \yes & 외 & 보류 & 경제활동 비율과 고용밀도 중 하나로 구체화. 소득과의 중복 확인.\\
인구밀도 & \maybe & \yes & \maybe & \yes & 외 & 보류 & 지역별 효과 차이. 무차량 가구ㆍ대중교통과 중복 확인.\\
토지이용 혼합ㆍ접근성 & \maybe & \maybe & \maybe & \yes & 외 & 보류 & 대표 지표 하나를 선택하고 구역 지표의 도시 집계 방법 확인.\\
대중교통ㆍ보행환경 & \maybe & \maybe & \maybe & \yes & 외 & 보류 & 대체ㆍ보완이 공존. 대안 경쟁력ㆍ무차량 가구와 중복 확인.\\
식당ㆍ여가시설 밀도 & + & \maybe & \yes & \yes & 외 & 보류 & SF 한 도시의 이용량 근거. 시설 정의ㆍ도시 전이 확인 후 검토.\\

%---------새로 추가한 후보------------
연령 구조 & \maybe & \yes & \yes & \yes & 외 & 선택 & 18--34세 인구 비율을 우선 검토. 젊은 층 이용 근거가 있으나 도시 수준 부호는 추후 확인.\\
학생 비율 & \maybe & \maybe & \maybe & \yes & 외 & 보류 & 연령과 중복. 캠퍼스ㆍ학생 요금지원 효과가 섞여 있음.\\
교육ㆍ기술 친화도 & \maybe & \maybe & \maybe & \yes & 혼 & 보류 & 학력과 기술 친화도를 구분. 직접 설문ㆍ근거 확인 필요.\\
기존 호출 이용습관 & \maybe & \maybe & \maybe & \yes & 혼 & 보류 & 결정 이전 일반 호출 이용자료 필요. 과거 이용량과 태도상 습관을 구분.\\
안전 인식ㆍ노출 & \maybe & \maybe & \maybe & \yes & 혼 & 보류 & 공통 안전 설문 필요. 선택 후 만족도와 사고ㆍ범죄 노출을 구분.\\
규제 허용 여부 & -- & \maybe & -- & -- & 별 & 분리 & 선택확률의 변수보다 운행 가능 조건으로 처리.\\
기후ㆍODD & \maybe & \maybe & \maybe & \yes & 혼 & 분리 & 기후는 외생. ODD와 운행 가용성은 기술ㆍ운영조건에 따라 판단.\\
서비스 개시 후 경과기간 & \maybe & \maybe & \yes & \yes & 내 & 보류 & 개시 시점이 운영 결정에 의존. 진입 후 시간 경로를 다루는 확장에서 검토.\\
\end{longtable}}

\section{1차 선택 결과: 외생ㆍ내생 표시}
\textbf{도시 특성 2개와 결정의존 서비스 변수 3개}를 우선 대상으로 정리한다. 내생 서비스 변수의 사용은 기준 8을 확인한 뒤 결정의존 ER-DRO 확장에서 반영한다. 아래의 외생ㆍ내생은 \emph{운영 결정에 의존하는지}에 따른 구분이다.

{\footnotesize
\begin{longtable}{@{}L{\fourwidth*\real{0.16}}L{\fourwidth*\real{0.30}}L{\fourwidth*\real{0.34}}L{\fourwidth*\real{0.20}}@{}}
\toprule 변수 & 근거 논문 & 논문에서 확인된 범위ㆍ결과 & 측정자료\\
\midrule\endfirsthead
\toprule 변수 & 근거 논문 & 논문에서 확인된 범위ㆍ결과 & 측정자료\\
\midrule\endhead
\bottomrule\endlastfoot
\textbf{무차량 가구 비율}\newline (외생) & \elk; \loa & 자가용 대안의 제약을 대표. 무차량 가구/전체 가구로 사용. & ACS B08201\\
\textbf{연령 구조}\newline (외생) & \loa; \rapid & 연령 구성을 대표. 18--34세 인구/총인구로 사용. 개인 선택에서의 부호는 일률적으로 고정하지 않음. & ACS B01001\\
\textbf{상대 요금}\newline (자사 요금: 내생) & \yin & 요금 차이에 대한 선택 반응. 자사 요금을 바꾸는 결정의존 확장용으로 선택. & 양 서비스 요금 + 선택결과\\
\textbf{상대 대기시간}\newline (자사 대기: 내생) & \yin & 대기시간 차이에 대한 선택 반응. 차량 수ㆍ배치와 대기의 관계를 연결해 사용. & 차량ㆍ배치ㆍ대기 기록 + 선택결과\\
\textbf{차내 지원기능}\newline (내생 후보) & \yin & 목적지 변경ㆍ운영자 연락 등의 선택효과. 기능을 선택하는 운영 확장용으로 유지. & 기능별 선택자료\\
\end{longtable}}

\textbf{추가 검토.} 소득 또는 경제활동 비율은 다음 추가 후보다. 변수 수는 관측 도시 수와 중복 여부를 보고 정한다. 경쟁서비스의 요금ㆍ대기는 외생 후보로 두되 관측 범위를 확인한다. 요금ㆍ기능이 현재 모형의 결정변수에 포함되지 않았다면, 표의 내생 표시는 확장 시의 역할을 뜻한다.

\textbf{별도 처리.} 규제는 운행 가능 조건, 기후ㆍODD는 운행 가용성으로 처리한다. 신뢰ㆍ지식ㆍ경험ㆍ안전 등은 공통 자료 확보 전 보류한다. 서비스 경과기간은 진입 후 확장에서 검토한다.

\textbf{사후 확인.} 자료를 모은 뒤 도시 전이(6), 잔차 패턴(7), 결정 반응(8)을 검증하여 최종 선택한다. $f(x)+Q(x)\epsilon$에서는 원래 잔차의 이분산을 무조건 탈락 사유로 삼지 않고, 표준화 잔차를 확인한다.



\section{아직 해야 할 일 / 한계}
\begin{itemize}
    \item p의 추정 위해 관측값(y) 도시 c의 로봇택시 통행량 데이터 필요
    
    \item \textbf{관측 p와 수요 D}: 완료 통행으로 계산한 \(p\)에는 공급 부족의 영향이 포함됨. 이를 그대로 수요로 쓰면 실현되지 못한 수요를 과소평가할 수 있습니다. 요청·거절 자료로 보정하거나, 공급 제약이 작다는 가정을 명시 필요

    \item p의 외생/내생변수들에 대한 데이터 추가 확인 필요
    
\end{itemize}

\end{document}
